\documentclass{ximera}
\title{Wedge products}

\begin{document}

    \begin{abstract} Wedge products, graded commutativity, and associativity. \end{abstract}
    \maketitle

    A reference for this material is Chapter 14 of John M. Lee. Introduction to smooth manifolds. Second edition. Grad. Texts in Math., 218. Springer, New York, 2013. xvi+708pp. ISBN: 978-1-4419-9981-8.

    \begin{definition}[Wedge product]
        Let $\omega \in \Omega ^k (M) $ and $\eta \in \Omega ^{\ell} (M)$, the \emph{wedge product} $\omega \wedge \eta \in \Omega ^{k + \ell } (M)$ is defined as
        \[     \omega \wedge \eta \, (X_1 , \ldots , X_{k + \ell} )  : = \frac{1}{\, k! \, \ell ! \, } \sum _{\sigma \in S_{k + \ell}} \text{sign}(\sigma)\, \omega (  X_{\sigma (1) }  , \ldots , X_{\sigma (k)})\,  \eta (  X_{\sigma (k+1)} , \ldots , X_{\sigma (k + \ell )}).                \]
    \end{definition}

    After renormalization, the wedge product is the antisymmetrization of the tensor product
        \[     \omega \wedge \eta = \frac{(k + \ell ) !}{k!\,  \ell ! } A_{k + \ell } (\omega \otimes \eta ) .                      \]
    \begin{proposition}[Graded commutativity]
        The wedge product is graded-commutative. That is, if $\omega \in \Omega ^k (M) $ and $\eta \in \Omega ^{\ell} (M)$, then
        \[   \omega \wedge \eta =  (-1) ^{k\ell} \, \eta \wedge \omega .           \]
    \end{proposition}

    \begin{solution}
    Let $\tau \in S_{k + \ell }$ denote the permutation 
    \[   1 \mapsto  k  + 1 , \ldots , \ell \mapsto k + \ell , \ell + 1 \mapsto 1 , \ldots ,  \ell + k \mapsto k .                         \]
    Note that $\text{sign}(\tau) = (-1)^{k \ell}$. Then 
    \begin{align*}
        k! \, \ell ! &  \, \omega  \wedge \eta  (X_1,  \ldots  , X_{k + \ell})  \\
        & = ( k + \ell ) !\, A_{k + \ell} (\omega \otimes \eta ) (X_1, \ldots , X_{k + \ell })\\
        & = \sum _{\sigma \in S_{k + \ell}} \text{sign}(\sigma)\, \omega (  X_{\sigma (1) }  , \ldots , X_{\sigma (k)})\,  \eta (  X_{\sigma (k+1)} , \ldots , X_{\sigma (k + \ell )})\\
        & = \sum _{\sigma \in S_{k + \ell}} \text{sign}(\sigma)\,   \eta (  X_{\sigma \tau (1 )} , \ldots , X_{\sigma \tau ( \ell )}) \, \omega (  X_{\sigma \tau  (\ell + 1) }  , \ldots , X_{\sigma \tau  (\ell + k)}) \\
        & = \text{sign}(\tau) \sum _{\sigma \in S_{k + \ell}} \text{sign}(\sigma)\,   \eta (  X_{\sigma  (1 )} , \ldots , X_{\sigma  ( \ell )}) \, \omega (  X_{\sigma  (\ell + 1) }  , \ldots , X_{\sigma   (\ell + k)}) \\
       & = (-1)^{k \ell}  k! \, \ell ! \, \eta \wedge \omega (X_1, \ldots , X_{k + \ell}).
    \end{align*}

    \end{solution}



\begin{proposition}[Associativity]
    The wedge product is associative. That is, if $\omega _j \in \Omega ^k_k (M)$ for $j \in \{ 1, 2, 3 \}$, then 
    \[     (\omega_1 \wedge \omega _2)  \wedge \omega _3 = \omega _1 \wedge (\omega _2 \wedge \omega _3).                      \]
\end{proposition}


\begin{solution} 
    \begin{align*}
           (\omega_1 \wedge \omega _2)  \wedge \omega _3 & = \dfrac{(k_1 + k_2)!}{k_1!\,k_2!} A_{k_1 + k_2} (\omega_1 \otimes \omega_2) \wedge A_{k_3} ( \omega_3 )  \\
           & = \dfrac{(k_1 + k_2 + k_3 )!}{k_1!\,k_2!\, k_3!} A_{k_1 + k_2 + k _3} ( A_{k_1 + k_2} (\omega_1 \otimes \omega_2) \otimes A_{k_3} ( \omega_3 )  )\\
           & = \dfrac{(k_1 + k_2 + k_3 )!}{k_1!\,k_2!\, k_3!} A_{k_1 + k_2 + k _3} ( \omega_1 \otimes \omega_2  \otimes  \omega_3  ).
    \end{align*}
    The last expression is associative, so the first one is. 
\end{solution}

\end{document}


